%------------------------------------------------------------------------------%
\section{Função modular}

Antes de definirmos função modular, precisamos do conceito de valor absoluto.

%------------------------------------------------------------------------------%
\begin{definicao}{}{def:}
  Dado um número real $x$, definimos o
  \emph{valor absoluto}\index{valor absoluto},
  ou \emph{módulo}\index{Módulo} de $x$ como
  \[
    \abs{a} = \begin{cases}
    \hpm x, & \text{ se } x \geq 0, \\
    -    x, & \text{ se } x <    0.
    \end{cases}
  \]
\end{definicao}
%------------------------------------------------------------------------------%

Note que a definição de valor absoluto é feita por partes. Se $x$ é um
número positivo, o valor absoluto é seu próprio valor.
Logo, $\abs{4}=4$.
Em contrapartida, se $x$ é um valor negativo, o valor absoluto é obtido
eliminando o sinal. Assim,
\[
  \abs{-10} = -(-10) = 10.
\]

%------------------------------------------------------------------------------%
\begin{example}
  Elimine o módulo da expressão
  \(
    \abs{-\frac{3}{2}}
  \)
  usando a definição de valor absoluto.

  \solution
  \begin{align*}
    \abs{-\frac{3}{2}}
    & = - \left(-\frac{3}{2} \right) \\
    & = \frac{3}{2}
  \end{align*}
\end{example}
%------------------------------------------------------------------------------%

%------------------------------------------------------------------------------%
\begin{example}
  Elimine o módulo da expressão
  \(
    -\abs{-5}
  \)
  usando a definição de valor absoluto.

  \solution
  Sabemos que $\abs{-5}=5$, então $-\abs{-5} =-5$.
\end{example}
%------------------------------------------------------------------------------%

%------------------------------------------------------------------------------%
\begin{example}
  Elimine o módulo da expressão
  \(
    \abs{3-2\sqrt{3}}
  \)
  usando a definição de valor absoluto.

  \solution
  Você deve atentar que
  \[
    3-2\sqrt{3}<0,
  \]
  logo
  \begin{align*}
   \abs{3-2\sqrt{3}}
   & = -\left( 3 - 2\sqrt{3} \right) \\
   & =  2\sqrt{3}-3.
  \end{align*}
\end{example}
%------------------------------------------------------------------------------%

%------------------------------------------------------------------------------%
\begin{example}
  Elimine o módulo da expressão
  \(
    \abs{2x+1}
  \)
  usando a definição de valor absoluto.

  \solution
  Nesse caso, usando a definição de módulo, temos
  \begin{align*}
    \abs{2x+1}
    & = \begin{cases}
          \hpm 2x+1,  & \text{ se } 2x + 1 \geq 0 \\
          -   (2x+1), & \text{ se } 2x + 1 <    0
        \end{cases}
    \\[2mm]
    & = \begin{cases}
        \hpm 2x+1, & \text{ se } x \geq -\sfrac{1}{2} \\
        -    2x-1, & \text{ se } x <    -\sfrac{1}{2}
        \end{cases}
  \end{align*}
\end{example}
%------------------------------------------------------------------------------%

A seguir seguem algumas propriedades do valor absoluto,
considerado $a$ e $b$ números reais.
\begin{gather}
  \abs{a} \geq 0                 \\[2mm]
  \abs{-a} = \abs{a}             \\[2mm]
  \abs{ab} = \abs{a} \; \abs{b}  \\[2mm]
  \abs{\frac{a}{b}} = \frac{\abs{a}}{\abs{b}} \qquad b \neq 0  \\[2mm]
  \abs{a^n} = a^n \qquad \text{ se } n \text{ é par}           \\[2mm]
  \abs{a+b} \leq \abs{a} + \abs{b}            \\[2mm]
  \abs{a-b} = a-b \qquad \text{ se } a \geq b \\[2mm]
  \abs{a-b} = b-a \qquad \text{ se } a \leq b \\[2mm]
  \sqrt{a^2} = \abs{a}
\end{gather}

Portanto,
\emph{função modular}\index{Função!modulas} ou
\emph{função valor absoluto}\index{Função!valor absolutos}
é uma função definida por partes, dada por
\[
  f(x)
  = \abs{x}
  = \begin{cases}
  \hpm x, & \text{ se } x \geq 0,\\
  -    x, & \text{ se } x <    0.
  \end{cases}
\]

%------------------------------------------------------------------------------%
\begin{figure}
  \centering
  \includegraphics[width=9cm,height=6cm]{modular}
  \caption{Gráfico da função modular.}
  \label{modulargr}
\end{figure}
%------------------------------------------------------------------------------%

A Figura~\ref{modulargr} mostra o gráfico da  função modular,
cujas características são:
\begin{itemize}
  \item O domínio é $\R$.
  \item O conjunto imagem é $[0, \infty)$.
  \item Há um zero em $x=0$.
  \item A função é decrescente em $(-\infty, 0)$
                 e   crescente em $(0, \infty)$.
  \item A função é par.
\end{itemize}

%------------------------------------------------------------------------------%
\begin{example}
  A partir do gráfico de $f(x)= \abs{x}$ faça o
  esboço do gráfico da função
  $f(x)= \abs{x-3} +4$

  \solution
  Nesse caso, para construir o gráfico de $f$ basta deslocar o gráfico
  \ref{modulargr} 3 unidades para a direita e 4 unidades para cima,
  conforme mostra a figura.

  \IncludeGraphicCentredEx{mod1}[0.5]
\end{example}
%------------------------------------------------------------------------------%

%------------------------------------------------------------------------------%
\begin{example}
  A partir do gráfico de $f(x)= \abs{x}$ faça o
  esboço do gráfico da função
  $f(x)= 2 \abs{x} -2$

  \solution
  Aqui, precisamos expandir o gráfico de $y=\abs{x}$ verticalmente por
  um fator 2 e descer 2 unidades, conforme mostra a figura.

  \IncludeGraphicCentredEx{mod2}[0.5]
\end{example}
%------------------------------------------------------------------------------%

Ao compormos uma função modular com uma função $f$ qualquer,
obtemos o módulo de $f$ que é dado por
\[
  \abs{f(x)} = \begin{cases}
  \hpm f(x), & \text{ se } f(x) \geq 0,\\
  -    f(x), & \text{ se } f(x) <    0.
  \end{cases}
\]

Em outras palavras, para fazer o esboço do gráfico da função
$\abs{f(x)}$, basta considerar os pedaços do gráfico de $f$
que estão abaixo do eixo $x$ e refleti-los com relação ao eixo $x$.

Por outro lado, se quisermos fazer o esboço do gráfico
\[
  f(x)= g(\abs{x}),
\]
ignoramos a parte do gráfico de $g(x)$ que está à esquerda do eixo $y$
e refletimos, com relação ao eixo $y$, a parte do gráfico de $g(x)$
que está à direita do eixo $y$ e então obtemos o gráfico da nova função.
Funções obtidas dessa maneira são sempre pares, já que $f(x) = f(-x)$.

%------------------------------------------------------------------------------%
\begin{example}
  Faça o esboço do gráfico da função
  \(
    f(x)= \abs{2x-9}
  \).

  \solution
  Nesse caso, temos que
  \[
    f(x) = g\left(h(x)\right),
    \quad\text{onde}\quad
    g(x) = \abs{x}
    \quad\text{e}\quad
    h(x) = 2x-9.
  \]
  Portanto, para obtermos o gráfico de $f$, basta construirmos o gráfico de
  $h(x)=2x-9$ e refletir o pedaço desse gráfico que está abaixo do eixo
  $x$,com respeito ao eixo $x$.

  $h(x)$ é uma função linear, cujo gráfico é
  uma reta, que aprendemos a esboçar no capítulo anterior.

  A figura mostra o gráfico de $f$.

  \IncludeGraphicCentredEx{mod3}[0.5]
\end{example}

%------------------------------------------------------------------------------%
\begin{example}
  Faça o esboço do gráfico da função
  \(
    g(x) = \abs{x^2-3x+2}
  \).

  \solution
  Do mesmo modo do item anterior, $g(x)$ é a composição de duas funções, ou
  seja,
  \[
    g(x) = p(q(x)),
    \quad\text{com}\quad
    p(x) = \abs{x}
    \quad\text{e}\quad
    q(x) = x^2 - 3x + 2.
  \]
  Logo, para construir o gráfico de $g$, basta fazermos o esboço
  do gráfico de $q$ e rebatermos a parte negativa com relação ao eixo $x$.

  Sabemos que $q(x)$ é uma função quadrática e seu gráfico é uma
  parábola com concavidade voltada para cima e o procedimento para
  o esboço do gráfico de $q$ foi apresentado no capítulo anterior.

  A figura mostra o gráfico de $g$.

  \IncludeGraphicCentredEx{mod4}[0.5]
\end{example}

%------------------------------------------------------------------------------%
\begin{example}
  Faça o esboço do gráfico da função
  \(
    f(x) = \abs{x}^2 - 4\abs{x}+3
  \).

  \solution
  Note que
  \[
    f(x) = g\left(\abs{x}\right),
    \quad\text{onde}\quad
    g(x) = x^2 - 4x + 3.
  \]
  Para fazer o esboço do gráfico de $f(x)$, ignoramos a parte do
  gráfico de $g(x)$ que está à esquerda do eixo $y$ e refletimos,
  com relação ao eixo $y$, a parte do gráfico de $g(x)$ que está
  a direita desse eixo.

  O gráfico de $g$ é uma parábola, com concavidade voltada para cima.

  Sendo assim, a figura apresenta o esboço do gráfico de $g$.

  \IncludeGraphicCentredEx{mod6}[0.5]
\end{example}
%------------------------------------------------------------------------------%
